\scanpage{47}
\mspage{11}
Par ailleurs $X_{\Delta_1}$ est le graphe d'une relation
d'ordre sur $X_{\Delta_0}=\{(x,i)\in X\times I\mid x\in X_i\}$
avec
\[
(8)\qquad (x,i)\leq(y,j)\ \Longleftrightarrow\ x=y,\quad i\leq j
\]
i.e. c'est la relation d'ordre induite par
l'ordre produit de $X\times I$, $X$ étant muni
de l'ordre discret. Par suite, le
couple $(X_{\Delta_0},X_{\Delta_1}\subset X_{\Delta_0}\times X_{\Delta_0})$ peut
être regardé comme un ``objet ordonné de
la catégorie des espaces topologiques'' ---
d'où une catégorie, dans la catégorie
des espaces topologiques, ayant $X_{\Delta_0}$ com-
me objet des objets, $\unsure{X}_{\Delta_1}$ comme objet
des flèches, $a$ et $b$ de (6) comme
morphismes source et but, et la
composition des flèches étant le
morphisme
\diagram{\begin{tikzcd}[column sep=large,row sep=small]
(9)\quad X_{\Delta_2}\overset{\mathrm{def}}{=}(X_{\Delta_1},b)\times_{X_{\Delta_0}}(X_{\Delta_1},a)\arrow[r]&X_{\Delta_1}\\
\cap&\cap\\
X_{\Delta_0}\times X_{\Delta_0}\times X_{\Delta_0}&X_{\Delta_0}\times X_{\Delta_0}
\end{tikzcd}}
induit par la projection $(x,y,z)\longrightarrow(x,z)$.
Plus généralement, on introduit \unsure{en|un}
\mspage{46}
